\documentclass[11pt,a4paper]{amsart}

%% ── Packages ──────────────────────────────────────────────────────────────
\usepackage[utf8]{inputenc}
\usepackage[T1]{fontenc}
\usepackage{amsmath,amssymb,amsthm,mathtools}
\usepackage{enumerate}
\usepackage{booktabs}
\usepackage{array}
\usepackage{hyperref}
\usepackage[capitalise,noabbrev]{cleveref}
\usepackage{microtype}
\usepackage{geometry}
\geometry{margin=1in}

%% ── Theorem environments ──────────────────────────────────────────────────
\theoremstyle{plain}
\newtheorem{theorem}{Theorem}[section]
\newtheorem{lemma}[theorem]{Lemma}
\newtheorem{proposition}[theorem]{Proposition}
\newtheorem{corollary}[theorem]{Corollary}
\newtheorem{conjecture}[theorem]{Conjecture}

\theoremstyle{definition}
\newtheorem{definition}[theorem]{Definition}
\newtheorem{observation}[theorem]{Observation}

\theoremstyle{remark}
\newtheorem{remark}[theorem]{Remark}

%% ── Custom commands ───────────────────────────────────────────────────────
\newcommand{\N}{\mathbb{N}}
\newcommand{\Z}{\mathbb{Z}}
\newcommand{\Q}{\mathbb{Q}}
\newcommand{\R}{\mathbb{R}}
\newcommand{\vp}{\varphi}
\newcommand{\myc}{\operatorname{Mycielski}}
\newcommand{\chr}{\chi}
\newcommand{\om}{\omega}
\newcommand{\Spec}{\operatorname{Spec}}
\newcommand{\spn}{\operatorname{span}}
\newcommand{\LHof}{L_{\mathrm{Hof}}}

%% ── Title block ───────────────────────────────────────────────────────────
\title[Golden Eigenvalues and the Tihany Conjecture for Mycielski Graphs]{%
  Golden Eigenvalues and the Erd\H{o}s--Lov\'asz Tihany Conjecture\\
  for Mycielski Graphs}

\author{V.\,F.\,S.\ Santos}
\date{February 2026}

\subjclass[2020]{05C15, 05C50, 05C31}
\keywords{Tihany conjecture, Mycielski graphs, golden ratio, spectral graph theory, chromatic number}

\begin{document}

\begin{abstract}
The Erd\H{o}s--Lov\'asz Tihany Conjecture (1968) asserts that every graph~$G$
with $\chr(G) \geq s+t-1 > \om(G)$ admits a vertex partition into parts with
chromatic numbers $\geq s$ and $\geq t$, respectively.  We prove the conjecture
for the infinite family of pairs $(3, k{-}2)$ on Mycielski graphs~$M_k$ for all
$k \geq 5$.

Our approach is spectral, centred on the golden ratio
$\vp = (1+\sqrt{5})/2$.  The pentagon~$C_5$---the minimal graph with
$\chr > \om$---has adjacency eigenvalues
$\{2, \vp^{-1}, \vp^{-1}, -\vp, -\vp\}$, and this golden spectral structure
propagates through the Mycielski construction: the golden ratio's defining
equation $\mu^2 - \mu - 1 = 0$ arises exactly from the Mycielski eigenvalue
relation (Lemma~\ref{lem:golden-propagation}).  We prove the \emph{$C_5$-Peeling
Existence Theorem}: for every $M_k$ ($k\geq 4$), a direction in the golden
eigenspace peels off a~$C_5$ from~$M_k$.  The proof is constructive via
\emph{spectral interferometry}: two Mycielski lift paths span a four-dimensional
subspace whose layer-control matrix has $\det = \sqrt{5}$, enabling independent
phase steering to select a \emph{diagonal lift}~$C_5$---the reverse cycle
through alternating address layers.  Computationally, the Hoffman margin of this
partition is $F = \vp^{-3} = \sqrt{5}-2$ exactly, verified for all $k \leq 12$.

The key advance is the \emph{Golden Sub-Induction}
(Theorem~\ref{thm:main}): for $k \geq 6$, the $1/\vp$ shadow attenuation forces the
peeled~$C_5$ into the original vertex block of~$M_k$, so the remainder
$M_k \setminus P_k$ contains $\myc(M_{k-1}\setminus P_{k-1})$ as a subgraph,
where $(P_k)_{k\geq 5}$ is a coherent family of peeled pentagons.
Since $\chr(\myc(G)) = \chr(G)+1$ for any graph with an edge, this yields the
inductive bound $\chr(M_k \setminus P_k) \geq k-2$.  Combined with
$\chr(C_5) = 3$, this settles the Tihany conjecture for the pair $(3,k{-}2)$ on
$M_k$ for every $k\geq 5$---an infinite family of previously open cases.
\end{abstract}

\maketitle

%% ══════════════════════════════════════════════════════════════════════════
\section{Introduction and context}\label{sec:intro}
%% ══════════════════════════════════════════════════════════════════════════

\subsection{The conjecture}

\begin{conjecture}[Erd\H{o}s--Lov\'asz Tihany Conjecture, 1968]
\label{conj:tihany}
For every graph~$G$ and integers $s,t \geq 2$ with
\[
  \chr(G) = s + t - 1 > \om(G),
\]
there exists a partition $V(G) = S \sqcup T$ such that
$\chr(G[S]) \geq s$ and $\chr(G[T]) \geq t$.
\end{conjecture}

The conjecture has been open since 1968.  Known cases are summarised in
\cref{tab:known-cases}.

\begin{table}[ht]
\centering
\caption{Known cases of the Erd\H{o}s--Lov\'asz Tihany Conjecture.}
\label{tab:known-cases}
\begin{tabular}{lll}
\toprule
Case & Method & Reference \\
\midrule
Small pairs $(s,t) \in \{(2,2),\ldots,(3,5)\}$
  & Direct / critical subgraph
  & \cite{stiebitz1987,sachs1993} \\
Line graphs
  & Structural (edge-colouring)
  & \cite{kostochka2003} \\
Quasi-line graphs
  & Structural
  & \cite{chudnovsky2008} \\
$\alpha(G) = 2$
  & Structural
  & \cite{balogh2009} \\
\textbf{$(3,t)$ for all $t\geq 3$ on Mycielski graphs}
  & \textbf{Spectral $C_5$-peeling + sub-induction}
  & \textbf{This paper (Theorem~\ref{thm:main})} \\
General & \textbf{OPEN} & \\
\bottomrule
\end{tabular}
\end{table}

\subsection{Main result}

\begin{theorem}[The Golden Sub-Induction]\label{thm:main}
For every integer $k \geq 5$, the Mycielski graph~$M_k$ satisfies the
Erd\H{o}s--Lov\'asz Tihany Conjecture for the pair
$(s,t) = (3, k{-}2)$.  That is, there exists a partition
$V(M_k) = S \sqcup T$ with $\chr(M_k[S]) \geq 3$ and
$\chr(M_k[T]) \geq k - 2$.
\end{theorem}

This settles the conjecture for the infinite family of pairs
$(3,3),(3,4),(3,5),(3,6),\ldots$ on the Mycielski graph class.  The pairs
$(3,t)$ for $t \geq 6$ (i.e., $k\geq 8$) are, to our knowledge, new.

The proof combines three ingredients, each fully proved in this paper:
\begin{enumerate}[(i)]
\item \emph{Spectral $C_5$-peeling existence} (Theorem~\ref{thm:c5-peeling}): a
  constructive proof that a golden eigenvector direction isolates a
  pentagon~$C_5$ from~$M_k$, via the non-singular scaling matrix of the
  Steering Theorem (Theorem~\ref{thm:steering}, $\det = \sqrt{5}$).
\item \emph{Coherent peeled family} (Lemma~\ref{lem:coherent-family}): for
  $k\geq 6$, the peeled~$C_5$ is the canonical image of the $(k{-}1)$-level
  pentagon in the original block.
\item \emph{Containment} (Lemma~\ref{lem:containment-coherent}): the remainder
  $M_k \setminus P_k$ contains $\myc(M_{k-1}\setminus P_{k-1})$ as a subgraph,
  enabling the inductive bound $\chr(M_k\setminus P_k)\geq k-2$.
\end{enumerate}
The base case $k = 5$ is established by an explicit certificate
(Lemma~\ref{lem:base-k5}: odd cycle + $3$-coloring of the remainder).  The Hoffman margin
characterisation $F = \vp^{-3}$ (Observation~\ref{obs:margin}) is a supplementary
computational result verified for $k \leq 12$; it is \emph{not} required for
the proof of Theorem~\ref{thm:main}.

\subsection{Why the golden ratio?}\label{ssec:why-golden}

Every graph with $\chr > \om$ is imperfect.  By the Strong Perfect Graph
Theorem~\cite{chudnovsky2006}, it contains an odd hole or odd antihole as an
induced subgraph.  The minimal odd hole is~$C_5$, whose adjacency spectrum is
\begin{equation}\label{eq:spec-c5}
  \Spec(C_5) = \bigl\{2,\;\vp^{-1},\;\vp^{-1},\;-\vp,\;-\vp\bigr\},
\end{equation}
where $\vp = (1+\sqrt{5})/2 \approx 1.618$.  This is exact: the eigenvalues of
$C_n$ are $2\cos(2\pi k/n)$, and $\cos(2\pi/5) = (\vp - 1)/2$.

The \emph{Lov\'asz theta function} evaluates to
\[
  \vartheta(C_5)
  = 1 - \frac{\lambda_{\max}}{\lambda_{\min}}
  = 1 + \frac{2}{\vp}
  = \sqrt{5}
  = \vp + \vp^{-1}.
\]
For~$C_5$ specifically, $\vartheta$ determines the Shannon capacity:
$\Theta(C_5) = \sqrt{5}$~\cite{lovasz1979}.  (This is a special property
of~$C_5$; for longer odd cycles $C_{2k+1}$ with $k\geq 3$, the Shannon
capacity remains unknown.)

The chromatic obstruction in~$C_5$ is algebraically governed by~$\vp$.  Our
central discovery is that this golden spectral structure \emph{propagates}
through the constructions that build high-chromatic graphs, and that
\emph{golden eigenvectors}---not the minimum eigenvector---provide natural
Tihany partitions.

%% ══════════════════════════════════════════════════════════════════════════
\section{Golden eigenvalue propagation}\label{sec:propagation}
%% ══════════════════════════════════════════════════════════════════════════

\subsection{The Mycielski construction}

The \emph{Mycielski functor}~$M$ takes a graph $G = (V,E)$ with
$V = \{v_1,\ldots,v_n\}$ and produces $M(G)$ with:
\begin{itemize}
\item \emph{Original vertices}: $v_1,\ldots,v_n$.
\item \emph{Shadow vertices}: $u_1,\ldots,u_n$, where each~$u_i$ is adjacent
  to $N(v_i)$, the neighbours of~$v_i$.
\item \emph{Apex vertex}: $w$, adjacent to all shadow vertices.
\end{itemize}
The key property is $\chr(M(G)) = \chr(G) + 1$ and $\om(M(G)) = \om(G)$.
Thus~$M$ preserves clique number but increments chromatic number---it is the
canonical factory for Tihany-relevant graphs~\cite{mycielski1955}.

\subsection{Computational evidence}

Starting from~$C_5$ with eigenvalues
$\{2,\vp^{-1},\vp^{-1},-\vp,-\vp\}$:

\begin{table}[ht]
\centering
\caption{Golden eigenvalues in iterated Mycielski graphs.}
\label{tab:golden-eigs}
\begin{tabular}{lcclc}
\toprule
Graph & $\chr$ & $\om$ & Golden eigenvalues (adjacency) & Multiplicity \\
\midrule
$C_5$             & 3 & 2 & $\vp^{-1}$ ($\times 2$), $-\vp$ ($\times 2$) & 4 of 5 \\
Gr\"otzsch $= M(C_5)$ & 4 & 2 & $-\vp^2$ ($\times 2$) & 2 of 11 \\
$M^2(C_5)$        & 5 & 2 & $\vp$ ($\times 7$), $-\vp^{-1}$ ($\times 7$) & 14 of 23 \\
$M^3(C_5)$        & 6 & 2 & $\vp^2$ ($\times 10$) & 10 of 47 \\
\bottomrule
\end{tabular}
\end{table}

The golden eigenvalue count \emph{grows} at each Mycielski iteration, with the
multiplicity increasing roughly as $n/2$.

\subsection{Statement and proof}

\begin{lemma}[Golden Propagation]\label{lem:golden-propagation}
For each eigenvalue~$\lambda$ of $A(G)$ whose eigenvector is orthogonal
to~$\mathbf{1}$, the Mycielski graph $M(G)$ has eigenvalues $\lambda\vp$ and
$-\lambda\vp^{-1}$.  In particular, if $\lambda \in \Q(\sqrt{5})$, then both
new eigenvalues lie in $\Q(\sqrt{5})$.
\end{lemma}

\begin{proof}
The adjacency matrix of $M(G)$ in the block decomposition (original / shadow /
apex) is
\[
  A\bigl(M(G)\bigr)
  = \begin{pmatrix}
      A       & A         & \mathbf{0} \\
      A       & O_n       & \mathbf{1} \\
      \mathbf{0}^T & \mathbf{1}^T & 0
    \end{pmatrix}.
\]
Let~$v$ be an eigenvector of~$A$ with eigenvalue~$\lambda$ and
$\mathbf{1}^T v = 0$.  Seek an eigenvector of $A(M(G))$ of the form
$(\beta v,\,\alpha v,\,0)^T$ with eigenvalue~$\mu$.  The eigenvalue equations
give
\[
  \lambda\beta + \lambda\alpha = \mu\beta,
  \qquad
  \lambda\beta = \mu\alpha.
\]
Eliminating $\alpha/\beta$ yields $\mu^2 - \lambda\mu - \lambda^2 = 0$.
Dividing by~$\lambda^2$:
\begin{equation}\label{eq:golden-char}
  \Bigl(\frac{\mu}{\lambda}\Bigr)^{\!2}
  - \frac{\mu}{\lambda} - 1 = 0.
\end{equation}
This is the \emph{defining equation of the golden ratio}.  The roots are
$\mu/\lambda = \vp$ and $\mu/\lambda = -\vp^{-1}$, giving
\[
  \mu = \lambda\vp
  \qquad\text{or}\qquad
  \mu = -\lambda\vp^{-1}.
\]
Since $\vp = (1+\sqrt{5})/2 \in \Q(\sqrt{5})$ and $\Q(\sqrt{5})$ is a field,
$\lambda\in\Q(\sqrt{5})$ implies both $\lambda\vp$ and $-\lambda\vp^{-1}$ lie
in $\Q(\sqrt{5})$.
\end{proof}

\begin{remark}\label{rem:remaining-eigs}
The remaining eigenvalues of $M(G)$ (from eigenvectors not orthogonal
to~$\mathbf{1}$, plus the apex-coupled modes) satisfy a different relation
involving the graph size~$n$.  For the $2(n-1)$ eigenvalues covered by
Lemma~\ref{lem:golden-propagation}, the golden scaling is exact (this is a theorem,
not a computation; the numerical check on $M^k(C_5)$ for $k = 1,2,3$ merely
confirms it).
\end{remark}

\subsection{Non-Mycielski examples}

The golden eigenvalue phenomenon extends beyond Mycielski:

\begin{table}[ht]
\centering
\caption{Golden eigenvalues in non-Mycielski graphs.}
\label{tab:non-mycielski}
\begin{tabular}{llc}
\toprule
Graph & Origin & Golden eigenvalues \\
\midrule
Icosahedron   & $H_3$ geometry & $\pm\sqrt{5}$ ($\times 3$ each) \\
Dodecahedron  & $H_3$ geometry & $\pm\sqrt{5}$ ($\times 3$ each) \\
Petersen $K(5,2)$ & Kneser & $\lambda_{\min}=-2$ (not golden, but $\chr_f = 5/2$) \\
\bottomrule
\end{tabular}
\end{table}

The icosahedron and dodecahedron have $\pm\sqrt{5} = \pm(\vp + \vp^{-1})$
eigenvalues because their symmetry group~$H_3$ is defined over $\Q(\sqrt{5})$.

%% ══════════════════════════════════════════════════════════════════════════
\section{Spectral certificate implies Tihany-valid cut}\label{sec:certificate}
%% ══════════════════════════════════════════════════════════════════════════

\subsection{Spectral lower bounds on \texorpdfstring{$\chr$}{χ}}

Two standard lower bounds on $\chr(H)$:

\medskip\noindent
\textbf{Hoffman bound}~\cite{hoffman1970}.  For any graph~$H$ with
$\lambda_{\min}(H) < 0$:
\begin{equation}\label{eq:hoffman}
  \LHof(H)
  \;:=\;
  1 - \frac{\lambda_{\max}(H)}{\lambda_{\min}(H)}
  \;\leq\; \chr(H).
\end{equation}

\noindent
\textbf{Lov\'asz theta}~\cite{lovasz1979} (sandwich theorem).
Writing $\overline{H}$ for the complement of~$H$:
\begin{equation}\label{eq:theta}
  \vartheta(\overline{H}) \leq \chr(H).
\end{equation}

\begin{remark}[Notation convention]\label{rem:theta-notation}
Throughout this paper, graph complement is always written with
$\overline{\,\cdot\,}$ (long bar): $\overline{H}$ denotes the complement
of~$H$ and $\overline{G[S]}$ the complement of the induced subgraph~$G[S]$.
By contrast, $\bar{\vartheta}$ denotes the Lov\'asz theta \emph{variant}
$\bar{\vartheta}(H) := \vartheta(\overline{H})$ and $\bar{\chr}$ denotes the
clique cover number; these are \emph{not} complements of functions.

The Lov\'asz sandwich theorem states
$\om(H) \leq \vartheta(\overline{H}) \leq \chr(H)$; equivalently,
$\alpha(H) \leq \vartheta(H) \leq \bar{\chr}(H)$.
The \textbf{key distinction} in this paper:
\begin{itemize}
\item The \emph{Hoffman bound} \eqref{eq:hoffman} acts on the graph~$H$
  \textbf{itself} (no complement);
\item The \emph{Lov\'asz theta certificate} \eqref{eq:theta-cert} applies
  $\vartheta$ to the \textbf{complement} $\overline{G[S]}$.
\end{itemize}
Both yield lower bounds on~$\chr(H)$, but the arguments are different.
\end{remark}

\subsection{Statement}

\begin{lemma}[Spectral certificate $\Rightarrow$ Tihany-valid cut]
\label{lem:spectral-certificate}
Let~$G$ satisfy $\chr(G) = s+t-1 > \om(G)$ with $s,t\geq 2$, and let
$V(G) = S\sqcup T$.  If either
\begin{equation}\label{eq:theta-cert}
  \vartheta\bigl(\overline{G[S]}\bigr) > s-1,
  \qquad
  \vartheta\bigl(\overline{G[T]}\bigr) > t-1,
\end{equation}
or
\begin{equation}\label{eq:hoffman-cert}
  1 - \frac{\lambda_{\max}(G[S])}{\lambda_{\min}(G[S])} > s-1,
  \qquad
  1 - \frac{\lambda_{\max}(G[T])}{\lambda_{\min}(G[T])} > t-1
\end{equation}
(with $\lambda_{\min}<0$ on each side), then
$\chr(G[S]) \geq s$ and $\chr(G[T]) \geq t$.
Hence $(S,T)$ is $(s,t)$-Tihany-valid.
\end{lemma}

\begin{proof}
The Lov\'asz sandwich theorem gives
$\vartheta(\overline{G[S]}) \leq \chr(G[S])$ (see Remark~\ref{rem:theta-notation}),
and the Hoffman bound gives
$1 - \lambda_{\max}(H)/\lambda_{\min}(H) \leq \chr(H)$ when
$\lambda_{\min}(H) < 0$.  So each strict inequality $> s-1$ (respectively
$> t-1$) implies $\chr \geq s$ (respectively $\chr \geq t$) by integrality
of~$\chr$.
\end{proof}

\begin{corollary}[Threshold-cut form]\label{cor:threshold}
Let~$\mathbf{x}$ be a unit eigenvector of $A(G)$, and for $\tau\in\R$ define
\[
  S_\tau = \{v : x_v \geq \tau\},
  \qquad
  T_\tau = V\setminus S_\tau.
\]
If there exists~$\tau$ such that the Hoffman pair (or the theta pair) of
Lemma~\ref{lem:spectral-certificate} holds for $(S_\tau, T_\tau)$, then~$G$
satisfies the Tihany conjecture for $(s,t)$.
\end{corollary}

\subsection{Finite optimisation formulation}\label{ssec:finite-opt}

Since $S_\tau$ only changes when~$\tau$ crosses a coordinate of~$\mathbf{x}$,
the function
\begin{equation}\label{eq:F-tau}
  F(\tau)
  := \min\!\Bigl\{
       \LHof\bigl(G[S_\tau]\bigr) - (s-1),\;
       \LHof\bigl(G[T_\tau]\bigr) - (t-1)
     \Bigr\}
\end{equation}
takes at most $|V(G)| + 1$ distinct values.  Therefore:

\begin{quote}
\itshape Existence of a Tihany-valid threshold cut is equivalent to
$\max_\tau F(\tau) \geq 0$.
\end{quote}

This gives a \emph{finite optimisation target}: the entire problem reduces to
showing $\max_\tau F(\tau) \geq 0$ for a specific eigenvector~$\mathbf{x}$ and
pair~$(s,t)$.

%% ══════════════════════════════════════════════════════════════════════════
\section{Computational evidence}\label{sec:evidence}
%% ══════════════════════════════════════════════════════════════════════════

\subsection{Key finding: golden eigenvectors, not \texorpdfstring{$\lambda_{\min}$}{λ\_min}}

We tested eigenvector-threshold cuts on all graphs with $\chr > \om$ in our
library.  The critical discovery:

\emph{The $\lambda_{\min}$-eigenvector often fails to give valid Tihany cuts.
But eigenvectors of golden eigenvalues succeed.}

\begin{table}[ht]
\centering
\caption{Golden eigenvector cuts vs.\ $\lambda_{\min}$ cuts.}
\label{tab:cuts}
\begin{tabular}{@{}lccccl@{}}
\toprule
Graph & $\chr$ & $\om$ & $\lambda_{\min}$ cut? & Golden cut? & Which~$\lambda$? \\
\midrule
$C_5$         & 3 & 2 & No  & Yes & $\vp^{-1}$ \\
Petersen      & 3 & 2 & Yes & Yes & $\lambda_{\min}=-2$ \\
Gr\"otzsch    & 4 & 2 & No  & Borderline & $-\vp^2$ (Hoffman tight) \\
$M_5$         & 5 & 2 & Borderline & \textbf{Yes (all 7 $\vp$-eigenvectors)} & $\vp$, $-\vp^{-1}$ \\
Icosahedron   & 4 & 3 & Yes & \textbf{Yes (margins $> 1.0$)} & $\pm\sqrt{5}$ \\
Dodecahedron  & 3 & 2 & Yes & Yes & $\pm\sqrt{5}$ \\
\bottomrule
\end{tabular}
\end{table}

\subsection{The icosahedron is spectacular}

The $\sqrt{5}$-eigenvector of the icosahedron at threshold $\tau = 0$ gives a
perfect $6/6$ split:
\begin{itemize}
\item $\LHof(G[S]) = 3.132$,
\item $\LHof(G[T]) = 3.132$,
\item both exceed $s-1 = 2$ by a margin of $1.13$.
\end{itemize}
The $H_3$ (icosahedral) symmetry, defined over $\Q(\sqrt{5})$, makes Tihany
almost trivially satisfied.

\subsection{Mycielski\texorpdfstring{$_5$}{₅}: the critical test}

$M_5$ ($\chr = 5$, $\om = 2$, $n = 23$) has $7$ eigenvalues equal to~$\vp$ and
$7$ equal to~$-\vp^{-1}$.

For the symmetric case $(s,t) = (3,3)$: \emph{every single} $\vp$-eigenvector
gives a valid Hoffman-certified Tihany cut.  Margins range from~$0.07$
to~$0.24$.

For the asymmetric cases $(2,4)$ and $(4,2)$: no eigenvector-threshold cut
found.  This is the main gap.

\subsection{Physical intuition}

Why do golden eigenvectors work where $\lambda_{\min}$ fails?
\begin{itemize}
\item The $\lambda_{\min}$-eigenvector encodes the \emph{strongest} spectral
  obstruction.  Cutting along it \emph{concentrates} the obstruction in one
  half.
\item Golden eigenvectors encode the \emph{pentagonal} obstruction---the
  algebraic content of the odd-hole structure.  Cutting along them \emph{splits
  the pentagon structure into both halves}.
\end{itemize}
Tihany requires chromatic complexity in \emph{both} parts.  Golden eigenvectors
distribute the pentagonal obstruction evenly; $\lambda_{\min}$ does not.

%% ══════════════════════════════════════════════════════════════════════════
\section{The \texorpdfstring{$C_5$}{C₅}-Peeling Theorem}\label{sec:peeling}
%% ══════════════════════════════════════════════════════════════════════════

\subsection{Summary of results}

\cref{tab:summary} distinguishes between \textbf{(A)}~fully proved results
and \textbf{(B)}~computationally supported claims.

\begin{table}[ht]
\centering
\caption{Summary of Tihany results for Mycielski graphs. Proof status:
\textbf{A}~= fully proved; \textbf{B}~= computationally verified.}
\label{tab:summary}
\begin{tabular}{@{}lllll@{}}
\toprule
Case & $(s,t)$ type & Method & Status & Layer \\
\midrule
Sub-critical: $s{+}t{-}1 < \chr$
  & any
  & Golden eigenvector + Hoffman
  & Proven
  & A \\
\textbf{Infinite family}
  & \textbf{$(3,k{-}2)$ for all $k\geq 5$}
  & \textbf{$C_5$-peeling + sub-induction}
  & \textbf{Proven (Theorem~\ref{thm:main})}
  & \textbf{A} \\
$C_5$-peeling + Hoffman margin
  & $(3,3)$ for $k\geq 5$
  & Spectral certification
  & Verified $k\leq 12$ (Observation~\ref{obs:margin})
  & B \\
Critical asymmetric: $(2,k{-}1)$
  & $k\geq 5$
  & Structural (vertex-criticality)
  & Partition exists (Proposition~\ref{prop:shielding})
  & A \\
\bottomrule
\end{tabular}
\end{table}

\subsection{Statement}\label{ssec:c5-peeling}

The main peeling result has two layers: a fully proved existence theorem and a
supplementary computational observation about the Hoffman margin.

\begin{theorem}[$C_5$-Peeling Existence]\label{thm:c5-peeling}
For every $k \geq 4$, the Mycielski graph~$M_k$ (with $\chr = k$, $\om = 2$,
$n_k$~vertices) admits a \emph{golden spectral partition}:

There exists a unit vector~$v$ in the $+\vp^{k-4}$ eigenspace of $A(M_k)$ and
a threshold~$\tau^*$ such that:
\begin{enumerate}[\upshape(1)]
\item The~$5$ vertices on the small side of the cut form a $5$-cycle:
  $G[T_{\tau^*}] \cong C_5$.
\item The Hoffman bound of the peeled pentagon is exact:
  $\LHof(C_5) = \sqrt{5} > 2$, certifying $\chr(C_5) \geq 3$.
\end{enumerate}
\end{theorem}

The proof of Theorem~\ref{thm:c5-peeling} is given in \cref{ssec:proof-peeling}, via
the Steering Theorem (Theorem~\ref{thm:steering}) and the diagonal-lift mechanism
(\cref{ssec:diagonal}).

\begin{observation}[Golden Margin]\label{obs:margin}
Computationally, for all $k = 5,\ldots,12$, the $C_5$-peeling cut of
Theorem~\ref{thm:c5-peeling} additionally satisfies:
\begin{enumerate}[\upshape(1)]
\setcounter{enumi}{2}
\item The Hoffman bound of the large part exceeds~$\sqrt{5}$:
  $\LHof(G[S_{\tau^*}]) > \sqrt{5}$ for all $k \geq 5$.
\item The margin is the third power of $\vp^{-1}$:
  \begin{equation}\label{eq:margin}
    F(\tau^*) = \vp^{-3} = \sqrt{5} - 2 \approx 0.2361.
  \end{equation}
\end{enumerate}
We conjecture that items (3) and (4) hold for all $k \geq 5$.  A complete
analytical proof would upgrade the spectral certificate from computational to
algebraic (see \cref{sec:open}).
\end{observation}

\begin{remark}\label{rem:margin-not-needed}
Theorem~\ref{thm:main} does \textbf{not} depend on items~(3)--(4) of
Observation~\ref{obs:margin}.  Its proof uses only the peeling existence
(Theorem~\ref{thm:c5-peeling}), the coherent peeled family
(Lemma~\ref{lem:coherent-family}), and containment
(Lemma~\ref{lem:containment-coherent}), with a direct computation for the base
case $k=5$.
The Hoffman margin is an independent structural characterisation of the
$C_5$-peeling partition.
\end{remark}

\subsection{Algebraic base case (\texorpdfstring{$k = 4$}{k=4}: Gr\"otzsch
graph)}\label{ssec:grotzsch}

\begin{proposition}[Gr\"otzsch $C_5$-peeling, exact]\label{prop:grotzsch}
Let $G = M(C_5)$ (the Gr\"otzsch graph, $M_4$ in our notation).  There exists
an eigenvector~$\mathbf{y}$ of $A(G)$ and a threshold~$\tau$ such that
$G[H_\tau] \cong C_5$.
\end{proposition}

\begin{proof}
Label~$C_5$ vertices $v_0,\ldots,v_4$ cyclically.  The vector
\[
  x_i = \cos\!\Bigl(\frac{4\pi i}{5}\Bigr),
  \qquad i = 0,\ldots,4,
\]
is an eigenvector of $A(C_5)$ with eigenvalue
$\lambda = 2\cos(4\pi/5) = -\vp$, and $\mathbf{1}^T\mathbf{x} = 0$.

By Lemma~\ref{lem:golden-propagation} (negative branch:
$\mu = -\lambda/\vp = 1$), the vector
\[
  \mathbf{y} = (\mathbf{x},\;-\vp\,\mathbf{x},\;0)^T
\]
is an eigenvector of $A(M(C_5))$ with eigenvalue $\mu = 1 = +\vp^0$.
Writing $a := \cos(2\pi/5) = (\sqrt{5}-1)/4 \approx 0.309$, the components
of~$\mathbf{y}$ are listed in \cref{tab:grotzsch-eig}.

\begin{table}[ht]
\centering
\caption{Eigenvector components for the Gr\"otzsch $C_5$-peeling.}
\label{tab:grotzsch-eig}
\begin{tabular}{lll}
\toprule
Vertex & Component of~$\mathbf{y}$ & Block \\
\midrule
$v_0$ & $+1$        & original \\
$v_1$ & $-\vp/2$    & original \\
$v_2$ & $+a$        & original \\
$v_3$ & $+a$        & original \\
$v_4$ & $-\vp/2$    & original \\
$u_0$ & $-\vp$      & shadow \\
$u_1$ & $+\vp^2/2$  & shadow \\
$u_2$ & $-1/2$      & shadow \\
$u_3$ & $-1/2$      & shadow \\
$u_4$ & $+\vp^2/2$  & shadow \\
$w$   & $0$         & apex \\
\bottomrule
\end{tabular}
\end{table}

For any threshold $\tau \in (0, a)$, the threshold set
$H_\tau = \{v : y_v \geq \tau\}$ consists of exactly~$5$ vertices:
\[
  H_\tau = \{v_0, v_2, v_3, u_1, u_4\}.
\]

Checking adjacencies in $M(C_5)$:
\begin{itemize}
\item Among base vertices: only $v_2 v_3$ is present (since $v_2, v_3$ are
  adjacent in~$C_5$).
\item $u_1$ (shadow of~$v_1$) is adjacent to $N(v_1) = \{v_0, v_2\}$.
\item $u_4$ (shadow of~$v_4$) is adjacent to $N(v_4) = \{v_0, v_3\}$.
\item No shadow--shadow edges.
\end{itemize}

So $G[H_\tau]$ has edge set
$\{v_0 u_1,\;u_1 v_2,\;v_2 v_3,\;v_3 u_4,\;u_4 v_0\}$, which is the
$5$-cycle $v_0 - u_1 - v_2 - v_3 - u_4 - v_0$.
\end{proof}

\begin{remark}\label{rem:grotzsch-1}
The eigenvalue $\mu = 1 = +\vp^0 = +\vp^{k-4}$ (with $k = 4$) matches the
universal pattern of Theorem~\ref{thm:c5-peeling}.  It arises from $C_5$'s eigenvalue
$-\vp$ via the \emph{negative} Mycielski branch ($\mu = -\lambda/\vp$), with
shadow/original ratio $-\vp$ (not $+1/\vp$).
\end{remark}

\begin{remark}\label{rem:grotzsch-2}
The $C_5$ eigenvector $\cos(4\pi i/5)$ belongs to the $m = 2$ Fourier mode
(eigenvalue $2\cos(4\pi/5) = -\vp$).  This explains the
$D_5$-representation structure in the Steering Theorem (Theorem~\ref{thm:steering}): the $m = 2,3$ irreps
dominate the peeling direction at all levels.
\end{remark}

\begin{remark}\label{rem:grotzsch-3}
At $k = 4$ the complement has $\LHof(G\setminus C_5) = 2.000$ exactly, giving
$F = 0$ (the borderline case).  For $k\geq 5$, the margin opens to
$F = \vp^{-3}$ (Observation~\ref{obs:margin}).
\end{remark}

\subsection{Computational verification}\label{ssec:comp-verify}

Theorem~\ref{thm:c5-peeling} and Observation~\ref{obs:margin} are verified exhaustively for
$k = 4$ through $k = 12$ by searching random directions in the $+\vp^{k-4}$
eigenspace; see \cref{tab:verification}.  Reproducibility details: for each~$k$,
we sample uniformly random unit vectors in the $+\vp^{k-4}$ eigenspace and
evaluate the threshold cut.  The ``Trial'' column reports the first successful
trial index (0-indexed).  Code is available at
\texttt{code/explore\_spectral\_tihany.py}.

\begin{table}[ht]
\centering
\caption{Exhaustive verification of $C_5$-peeling for $M_k$, $k = 4,\ldots,12$.}
\label{tab:verification}
\begin{tabular}{@{}rrcrccc@{}}
\toprule
$k$ & $n$ & Eigenvalue & Mult.\ & Trial & $\LHof(G[S])$ & $F(\tau^*)$ \\
\midrule
 4 &   11 & $+\vp^0 = 1$       &  5 &   2 & 2.0000 & 0.0000 \\
 5 &   23 & $+\vp^1 = 1.618$   &  7 &   8 & 2.2855 & \textbf{0.2361} \\
 6 &   47 & $+\vp^2 = 2.618$   &  9 &  71 & 2.3891 & \textbf{0.2361} \\
 7 &   95 & $+\vp^3 = 4.236$   & 11 &   0 & 2.5006 & \textbf{0.2361} \\
 8 &  191 & $+\vp^4 = 6.854$   & 13 &  55 & 2.5623 & \textbf{0.2361} \\
 9 &  383 & $+\vp^5 = 11.090$  & 15 &  37 & 2.6035 & \textbf{0.2361} \\
10 &  767 & $+\vp^6 = 17.944$  & 17 & 241 & 2.6253 & \textbf{0.2361} \\
11 & 1535 & $+\vp^7 = 29.034$  & 19 & 134 & 2.6433 & \textbf{0.2361} \\
12 & 3071 & $+\vp^8 = 46.979$  & 21 & 499 & 2.6550 & \textbf{0.2361} \\
\bottomrule
\end{tabular}
\end{table}

Key observations:
\begin{enumerate}[(1)]
\item $F = \vp^{-3}$ \emph{exactly} for all $k \geq 5$.  The bottleneck is
  always the $C_5$ part ($\LHof(C_5) = \sqrt{5} = 2 + \vp^{-3}$); the large
  part always exceeds~$\sqrt{5}$.
\item $\LHof(G[S])$ increases monotonically toward $\LHof(M_k)$, since removing
  $5$~vertices from an exponentially growing graph has vanishing effect.
\item $\LHof(M_k)$ converges to $\approx 2.66$: both $\lambda_{\max}$ and
  $|\lambda_{\min}|$ of~$M_k$ grow at rate~$\vp$ per level, and their ratio
  stabilises.
\item Eigenspace search is essential: individual eigenvectors fail at
  $k = 6, 9, 10$ because the $C_5$-peeling direction is a non-trivial linear
  combination in the eigenspace.  The multiplicity grows from~$5$ to~$21$, and
  the peeling direction becomes increasingly ``hidden.''
\end{enumerate}

\subsection{The diagonal-lift mechanism}\label{ssec:diagonal}

The peeled~$C_5$ is \emph{not} the original $C_5$ subgraph of~$M_k$.  It is a
\emph{diagonal lift} through the Mycielski tower: one descendant of each
original $C_5$ vertex, drawn from different address layers.

\medskip\noindent
\textbf{Address notation.}
Each non-apex vertex of~$M_k$ has a binary address
$(\sigma_{k-3},\ldots,\sigma_1) \in \{O,S\}^{k-3}$, where $\sigma_i = O$
(original) or $S$ (shadow) at Mycielski step~$i$ (step~$1$ denotes
$C_5 \to\text{Gr\"otzsch}$, step~$2$ denotes
$\text{Gr\"otzsch}\to M_5$, etc.), plus a base $C_5$ index
$j \in \{0,1,2,3,4\}$.

\medskip\noindent
\textbf{The reverse-cycle pattern (empirical).}
The following concrete description of the peeled~$C_5$ is an empirical
observation (verified for $k = 5$ through $k = 10$); it is \emph{not} required
by the proof of Theorem~\ref{thm:main}, which uses only the \emph{existence} of a
peeling direction (Theorem~\ref{thm:steering}).  In every case tested, the peeled~$C_5$
consists of the same~$5$ vertex numbers $(v_3, v_6, v_9, v_{11}, v_{13})$
forming the cycle
\[
  v_3 \to v_9 \to v_{13} \to v_6 \to v_{11} \to v_3
\]
with $C_5$ positions $4 \to 3 \to 2 \to 1 \to 0 \to 4$---the \emph{reverse
cycle} (each step decrements by~$1\bmod 5$).  The vertices come from exactly
three address types; see \cref{tab:address}.

\begin{table}[ht]
\centering
\caption{Address structure of the peeled diagonal $C_5$.}
\label{tab:address}
\begin{tabular}{llcl}
\toprule
Vertex & Address & $C_5$ position & Address type \\
\midrule
$v_3$  & $\underbrace{O\cdots O}_{k-3}.4$ & 4 & OO (all-original) \\
$v_9$  & $S\underbrace{O\cdots O}_{k-4}.3$ & 3 & SO \\
$v_{13}$ & $OS\underbrace{O\cdots O}_{k-5}.2$ & 2 & OSO \\
$v_6$  & $S\underbrace{O\cdots O}_{k-4}.1$ & 1 & SO \\
$v_{11}$ & $OS\underbrace{O\cdots O}_{k-5}.0$ & 0 & OSO \\
\bottomrule
\end{tabular}
\end{table}

The two SO vertices (positions $3$, $1$) and two OSO vertices (positions $2$,
$0$) alternate around the cycle, with the single OO vertex (position~$4$)
bridging them.  This pattern is universal: the same vertices peel at every
$k \geq 5$, with addresses growing by prepending~$O$'s.

\medskip\noindent
\textbf{Why the diagonal lift forms a~$C_5$.}
The Mycielski shadow adjacency rules guarantee it.  Each edge type is verified
in \cref{tab:edge-verify}.

\begin{table}[ht]
\centering
\caption{Edge verification of the diagonal-lift $C_5$.}
\label{tab:edge-verify}
\begin{tabular}{lll}
\toprule
Edge & Type & Rule \\
\midrule
$v_3 \leftrightarrow v_9$ & OO$\leftrightarrow$SO
  & shadow($C_5$ vtx 4) connects to $N_{C_5}(4) \ni 3$ \\
$v_9 \leftrightarrow v_{13}$ & SO$\leftrightarrow$OSO
  & shadow(Gr\"otzsch vtx 2) connects to $N_{\text{Gr\"otzsch}}(2) \ni 9$ \\
$v_{13} \leftrightarrow v_6$ & OSO$\leftrightarrow$SO
  & shadow(Gr\"otzsch vtx 2) connects to $N_{\text{Gr\"otzsch}}(2) \ni 6$ \\
$v_6 \leftrightarrow v_{11}$ & SO$\leftrightarrow$OSO
  & shadow(Gr\"otzsch vtx 0) connects to $N_{\text{Gr\"otzsch}}(0) \ni 6$ \\
$v_{11} \leftrightarrow v_3$ & OSO$\leftrightarrow$OO
  & shadow(Gr\"otzsch vtx 0) connects to $N_{\text{Gr\"otzsch}}(0) \ni 3$ \\
\bottomrule
\end{tabular}
\end{table}

Only the first two Mycielski steps matter for adjacency: the cross-layer edges
are determined by $N_{C_5}$ and $N_{\text{Gr\"otzsch}}$, and all subsequent
steps simply preserve these connections in the original block while extending
them into new shadow blocks.

\subsection{Spectral growth and Hoffman convergence}\label{ssec:spectral-growth}

Both $\lambda_{\max}(M_k)$ and $|\lambda_{\min}(M_k)|$ grow at rate~$\vp$ per
Mycielski step; see \cref{tab:spectral-growth}.

\begin{table}[ht]
\centering
\caption{Spectral growth of Mycielski graphs.}
\label{tab:spectral-growth}
\begin{tabular}{rcccc}
\toprule
$k$ & $\lambda_{\max}$ & $|\lambda_{\min}|$
  & Growth of $\lambda_{\max}$ & Growth of $|\lambda_{\min}|$ \\
\midrule
 3 &   2.000 &   1.618 & ---   & --- \\
 4 &   3.702 &   2.702 & 1.851 & 1.670 \\
 5 &   6.525 &   4.437 & 1.763 & 1.642 \\
 8 &  31.189 &  19.375 & $\to\vp$ & $\to\vp$ \\
12 & 224.015 & 135.161 & 1.629 & 1.623 \\
\bottomrule
\end{tabular}
\end{table}

Both growth rates converge to~$\vp$ from above, so the Hoffman bound
$\LHof(M_k) = 1 + \lambda_{\max}/|\lambda_{\min}|$ converges to a constant
$\approx 2.66$, well above $\sqrt{5} \approx 2.236$.  Since removing
$5$~vertices from a graph with $n_k \to \infty$ vertices has negligible
spectral effect, $\LHof(G[S]) > \sqrt{5}$ holds for all $k \geq 5$.

\begin{observation}[Large-part Hoffman bound: computed range]\label{obs:hoffman-large}
For $k = 5,\ldots,12$, the peeled remainder satisfies
$\LHof(M_k \setminus P_k) > \sqrt{5}$, with the bound increasing
monotonically toward $\LHof(M_k)$.
\end{observation}

\begin{conjecture}[Uniform large-part Hoffman gap]\label{conj:hoffman-large}
For all $k \geq 5$, $\LHof(M_k \setminus P_k) > \sqrt{5}$.
\end{conjecture}

\subsection{Proof of the \texorpdfstring{$C_5$}{C₅}-Peeling Theorem}
\label{ssec:proof-peeling}

The complete proof of Theorem~\ref{thm:c5-peeling} decomposes into three parts.

\medskip\noindent
\textbf{Part~A (Eigenspace existence).}
The eigenvalue $+\vp^{k-4}$ arises from $C_5$'s eigenvalue $-\vp$ by one
negative branch ($-\vp \to +1$) and $(k-4)$ positive branches
($+1 \to +\vp \to \cdots \to +\vp^{k-4}$).  (Equivalently, from $C_5$'s
eigenvalue $\vp^{-1}$ by $(k-3)$ positive branches.)  The eigenspace has
multiplicity $\geq 2$ (inherited from $C_5$'s two-dimensional eigenspaces).  By
Proposition~\ref{prop:grotzsch}, the Mycielski lift of the Gr\"otzsch
eigenvector~$\mathbf{y}$ lies exactly in this eigenspace at every subsequent
level.

\medskip\noindent
\textbf{Part~B ($C_5$-peeling direction).}
Although the golden eigenspace has multiplicity $2k{-}3$ (growing linearly), the
$C_5$-peeling direction always lies in a \emph{fixed four-dimensional subspace}
$\mathcal{S}_k$, reducing the problem from high-dimensional asymptotics to
finite, tractable geometry.

\begin{theorem}[Steering]\label{thm:steering}
For every $k \geq 5$, the $+\vp^{k-4}$ eigenspace of~$M_k$ contains a
four-dimensional subspace~$\mathcal{S}_k$ and a unit vector~$v \in
\mathcal{S}_k$ whose threshold cut peels a diagonal-lift $C_5$.  Concretely:
\begin{enumerate}[\upshape(i)]
\item \emph{Construction.}  $\mathcal{S}_k =
  \spn\!\bigl\{z_A^{(\cos)}, z_A^{(\sin)},
  z_B^{(\cos)}, z_B^{(\sin)}\bigr\}$,
  where $z_P^{(m)}$ is the Mycielski lift of mode~$m \in \{\cos,\sin\}$ along
  path~$P \in \{A,B\}$ (defined below).
\item \emph{Surjectivity.}  The amplitude mapping
  $\mathcal{S}_k \to \R^2$ that sends $v \in \mathcal{S}_k$ to the phase-layer
  amplitudes at the SO and OS address layers is surjective, governed by
  the invertible matrix
  \[
    M = \begin{pmatrix} -\vp & \vp^{-1} \\ \vp^{-1} & -\vp \end{pmatrix},
    \qquad \det(M) = \sqrt{5} \neq 0.
  \]
\item \emph{Existence.}  There exists a path-mixing vector
  $(e_A, e_B) = M^{-1}(\text{target}_{\mathrm{SO}},
  \text{target}_{\mathrm{OS}})^T$ and a mode phase $\theta$ such that
  $v = e_A z_A^{(\theta)} + e_B z_B^{(\theta)} \in \mathcal{S}_k$
  selects the diagonal-lift $C_5$.
\end{enumerate}
\end{theorem}

\begin{proof}
We establish each item in order.

\medskip
\noindent\textbf{Step~1: Construction and eigenspace membership.}
The eigenvalue $-\vp$ of~$C_5$ has a two-dimensional eigenspace spanned by:
\[
  x_j^{(\cos)} = \cos(4\pi j/5),
  \qquad
  x_j^{(\sin)} = \sin(4\pi j/5),
  \qquad j = 0,\ldots,4.
\]
Each mode lifts from $-\vp$ to $+\vp^{k-4}$ through $(k-3)$ Mycielski steps
via two distinct paths:
\begin{itemize}
\item \textbf{Path~A} (neg-first): negative branch at step~$1$
  ($-\vp \to +1$), then $(k-4)$ positive branches
  ($+1 \to +\vp \to \cdots \to +\vp^{k-4}$).
\item \textbf{Path~B} (pos-first): positive branch at step~$1$
  ($-\vp \to +\vp^2$), negative branch at step~$2$
  ($+\vp^2 \to +\vp$), then $(k-4)$ positive branches
  ($+\vp \to \cdots \to +\vp^{k-4}$).
\end{itemize}
Both paths terminate at $+\vp^{k-4}$.  Each step applies
Lemma~\ref{lem:golden-propagation}, so the resulting vectors
$z_A^{(\cos)}, z_A^{(\sin)}, z_B^{(\cos)}, z_B^{(\sin)}$ are
$+\vp^{k-4}$-eigenvectors by construction.

\medskip
\noindent\textbf{Step~2: Linear independence ($\dim \mathcal{S}_k = 4$).}
The two $C_5$ modes are orthogonal ($\langle x^{(\cos)}, x^{(\sin)}\rangle
= 0$), and this orthogonality is preserved through each Mycielski lift since
the lift $(\beta v, \alpha v, 0)^T$ is a scalar multiple of the original
pattern on each address layer.  It remains to show that paths A and B yield
linearly independent vectors.

By Lemma~\ref{lem:golden-propagation}, each Mycielski lift multiplies the
amplitude by a factor that depends on the branch chosen:
\begin{itemize}
\item A negative branch at eigenvalue~$\lambda$ multiplies the original-block
  amplitude by $\beta = \lambda/\mu$ and the shadow-block amplitude by
  $\alpha = 1$ (relative to the original).
\item A positive branch multiplies differently: $\alpha/\beta = 1/\vp$
  (positive branch) or $-\vp$ (negative branch).
\end{itemize}
Propagating through the first two Mycielski steps, the scaling factors at
each address prefix are shown in \cref{tab:scaling}.

\begin{table}[ht]
\centering
\caption{Address-layer scaling factors for Paths~A and~B.}
\label{tab:scaling}
\begin{tabular}{lccc}
\toprule
Address prefix & Path~A factor & Path~B factor & Status \\
\midrule
OO (both original) & $+1$        & $+1$        & Locked (identical) \\
SO (shadow step~1)  & $-\vp$     & $+\vp^{-1}$ & Free (distinct) \\
OS (shadow step~2)  & $+\vp^{-1}$ & $-\vp$     & Free (distinct) \\
SS (both shadow)   & $-1$        & $-1$        & Locked (identical) \\
\bottomrule
\end{tabular}
\end{table}

\noindent
Since $-\vp \neq +\vp^{-1}$ (they differ in both sign and magnitude), the
restriction of~$z_A^{(m)}$ and~$z_B^{(m)}$ to the SO layer are not scalar
multiples of each other.  Hence $z_A^{(m)}$ and $z_B^{(m)}$ are linearly
independent for each mode~$m$.  Combined with the $\cos/\sin$ orthogonality,
$\dim \mathcal{S}_k = 4$.

\medskip
\noindent\textbf{Step~3: Surjectivity of the layer map.}
Define the \emph{layer map} $\Phi \colon \mathcal{S}_k \to \R^2$ by
$\Phi(v) = (\text{SO-amplitude}(v),\, \text{OS-amplitude}(v))$.
On the path-mixing plane (fixing a $C_5$ mode), the map acts as
\[
  \Phi(e_A z_A^{(m)} + e_B z_B^{(m)})
  = M \begin{pmatrix} e_A \\ e_B \end{pmatrix},
  \qquad
  M = \begin{pmatrix} -\vp & \vp^{-1} \\ \vp^{-1} & -\vp \end{pmatrix}.
\]
We compute the determinant:
\begin{equation}\label{eq:det-sqrt5}
  \boxed{\det(M) = (-\vp)(-\vp) - (\vp^{-1})(\vp^{-1})
  = \vp^2 - \vp^{-2} = \vp + \vp^{-1} = \sqrt{5} \neq 0.}
\end{equation}
The identity $\vp^2 - \vp^{-2} = \sqrt{5}$ follows from
$\vp = (1+\sqrt{5})/2$: we have $\vp^2 = (3+\sqrt{5})/2$ and
$\vp^{-2} = (3-\sqrt{5})/2$, so their difference is $\sqrt{5}$.
Since $\det(M) \neq 0$, the map~$\Phi$ is surjective.

\medskip
\noindent\textbf{Step~4: Existence of the diagonal-lift direction.}
The diagonal-lift~$C_5$ requires the pentagonal phase at the SO layer to be
offset by $\Delta\theta = 2\pi/5$ from the OS layer (see
\cref{ssec:diagonal}).  Set
$(\text{target}_{\mathrm{SO}}, \text{target}_{\mathrm{OS}})^T$ to be the
desired amplitude pair.  Since~$M$ is invertible:
\[
  \begin{pmatrix} e_A \\ e_B \end{pmatrix}
  = M^{-1} \begin{pmatrix} \text{target}_{\mathrm{SO}} \\
  \text{target}_{\mathrm{OS}} \end{pmatrix},
  \qquad
  M^{-1} = \frac{1}{\sqrt{5}}
  \begin{pmatrix} -\vp & -\vp^{-1} \\ -\vp^{-1} & -\vp \end{pmatrix}.
\]
This yields an explicit $v \in \mathcal{S}_k$ whose threshold cut selects the
diagonal-lift~$C_5$.

\medskip
\noindent\textbf{$k$-independence.}
The scaling matrix~$M$ depends only on the \emph{first two} Mycielski steps
(where the paths diverge).  All subsequent steps are positive branches for
both paths, contributing \emph{identical} multiplicative factors that cancel
in the ratio.  Hence $M$ is the same at every~$k \geq 5$, and
$\det(M) = \sqrt{5}$ is $k$-independent.
\end{proof}

\begin{remark}
The proof gives an explicit linear-algebraic construction; no numerical search
is needed.  The computational verification (\cref{tab:verification}) confirms
that the constructed direction produces the expected $C_5$-peeling at every
tested~$k$, but the proof itself is purely algebraic.
\end{remark}

\medskip\noindent
\textbf{Part~C (Peeled subgraph is~$C_5$).}
The diagonal-lift adjacency verification (\cref{tab:edge-verify}) establishes
that $G[T_{\tau^*}] \cong C_5$ for the~$5$ vertices selected by the steering
direction.  Since $\LHof(C_5) = \sqrt{5} > 2$, we have $\chr(C_5) \geq 3$
(which also follows trivially from the odd-cycle structure).  This completes
the proof of Theorem~\ref{thm:c5-peeling}. \qed

\subsection{Eigenvector block structure}

For the golden eigenvalues of $M(G)$ arising from an eigenvalue~$\lambda$
of~$G$ with eigenvector $v \perp \mathbf{1}$:

\begin{table}[ht]
\centering
\caption{Block structure of golden eigenvectors.}
\label{tab:block}
\begin{tabular}{lll}
\toprule
Eigenvalue of $M(G)$ & Eigenvector of $M(G)$ & Shadow/Original ratio \\
\midrule
$\mu = \lambda\vp$        & $(\beta v,\;\alpha v,\;0)$   & $\alpha/\beta = 1/\vp$ \\
$\mu = -\lambda/\vp$      & $(\beta' v,\;\alpha' v,\;0)$ & $\alpha'/\beta' = -\vp$ \\
\bottomrule
\end{tabular}
\end{table}

These ratios follow directly from the proof of
Lemma~\ref{lem:golden-propagation}: the ansatz $(\beta v,\,\alpha v,\,0)^T$ yields
$\alpha/\beta = \lambda/\mu$, which evaluates to $1/\vp$ for the positive
branch and $-\vp$ for the negative branch.  (Numerical confirmation on the
Gr\"otzsch graph: agreement to machine precision, $\sigma < 10^{-6}$.)
However, the $C_5$-peeling direction in the full eigenspace is a
\emph{non-trivial combination} of these block-structured vectors.

\subsection{Schur complement field preservation}

\begin{proposition}\label{prop:schur}
If the graph Laplacian $L \in M_n(\Q(\sqrt{5}))$ and the partition $V = S \sqcup F$
has $L_{FF}$ invertible, then the Kron reduction
$L_{\mathrm{eff}} = L_{SS} - L_{SF}\,L_{FF}^{-1}\,L_{FS}$ satisfies
$L_{\mathrm{eff}} \in M_{|S|}(\Q(\sqrt{5}))$.
\end{proposition}

\begin{proof}
$\Q(\sqrt{5})$ is a field; the Schur complement uses only addition,
multiplication, and inversion of entries, all of which are closed in a field.
\end{proof}

\begin{remark}\label{rem:schur-caveat}
Kron reduction produces the \emph{effective Laplacian} (with effective edges),
not the Laplacian of the induced subgraph~$G[S]$.  For Hoffman bound
certification, we need the induced subgraph.  The Schur complement thus
preserves algebraic structure but does not directly control the Hoffman bound.
\end{remark}

\subsection{The asymmetric case: vertex-criticality argument}
\label{ssec:asymmetric}

For asymmetric critical pairs like $(2, \chr{-}1)$, $C_5$-peeling gives
$\chr \geq 3$ on both sides, which is insufficient when $s$ or $t > 3$.
However, the partition \emph{exists structurally}.

\begin{theorem}[Mycielski vertex-criticality; {\cite{mycielski1955}}]
\label{thm:vertex-crit}
Mycielski graphs~$M_k$ are $k$-vertex-critical:
$\chr(M_k - v) = k - 1$ for every vertex~$v$.
\end{theorem}

\begin{proposition}[Conditional asymmetric partition]\label{prop:asym-cond}
If an edge $uv \in E(M_k)$ satisfies $\chr(M_k - \{u,v\}) \geq k - 1$, then
$S = \{u,v\}$, $T = V(M_k) \setminus \{u,v\}$ is a $(2, k{-}1)$-Tihany
partition.
\end{proposition}

\begin{proof}
We have $\chr(M_k[S]) = \chr(K_2) = 2 \geq 2$ and
$\chr(M_k[T]) \geq k-1$ by hypothesis.  The Tihany condition
$\chr(M_k) = k \geq 2 + (k{-}1) - 1 = k$ and
$\om(M_k) = 2 < k$ are satisfied for all $k \geq 3$.
\end{proof}

\begin{observation}\label{obs:edge-verify}
Exhaustive verification on~$M_5$: for all~$71$ edges $\{u,v\}$,
$\chr(M_5 - \{u,v\}) = 4 = k - 1$.  Hence every edge of~$M_5$ yields a
$(2,4)$-Tihany partition via Proposition~\ref{prop:asym-cond}.
\end{observation}

\begin{remark}
Note that $k$-vertex-criticality (Theorem~\ref{thm:vertex-crit}) guarantees
$\chr(M_k - v) = k-1$ for removing a \emph{single} vertex, but does not by
itself imply $\chr(M_k - \{u,v\}) \geq k-1$ for an edge $\{u,v\}$: removing
two adjacent vertices could in principle drop $\chr$ by~$2$.  Extending
Observation~\ref{obs:edge-verify} to all~$k$ is Open Problem~2.
\end{remark}

\medskip\noindent
\textbf{Certification gap and the shielding effect.}
The Hoffman bound of $M_5 - \{u,v\}$ is at most~$2.41$, far below the needed
threshold of~$3$.  Moreover, this gap is \emph{fundamental}:

\begin{proposition}[Spectral Shielding]\label{prop:shielding}
The three standard spectral and convex-relaxation bounds---the Hoffman bound
$\LHof$, the Lov\'asz theta function $\bar{\vartheta}$, and the fractional
chromatic number $\chr_f$---all satisfy the strict inequality
$\LHof(M_k) \leq \bar{\vartheta}(M_k) \leq \chr_f(M_k) < 4$ for every
$k \geq 3$.  In particular, none of these bounds can certify $\chr \geq 4$
for any Mycielski graph.
\end{proposition}

\begin{proof}
The Larsen--Propp--Ullman recurrence~\cite{larsen1995} gives
$\chr_f(M(G)) = \chr_f(G) + 1/\alpha(G)$.  Since $\alpha(M_k) = n_{k-1}$
(the shadow vertices form a maximum independent set), the fractional chromatic
number satisfies
\[
  \chr_f(M_k)
  = \frac{5}{2} + \frac{1}{2} + \frac{1}{5} + \frac{1}{11}
    + \frac{1}{23} + \cdots + \frac{1}{n_{k-2}},
\]
where $n_j = 3\cdot 2^{j-2} - 1$.  This is a convergent series:
\[
  \chr_f(M_\infty)
  = \frac{5}{2} + \sum_{j=3}^{\infty} \frac{1}{\alpha(M_j)}
  \approx 3.377 < 4.
\]
By the Lov\'asz sandwich theorem,
$\om(G) \leq \bar{\vartheta}(G) \leq \chr_f(G) \leq \chr(G)$.  Therefore
\[
  \bar{\vartheta}(M_k) \leq \chr_f(M_k) < 4 \quad\text{for all } k.
\]
Since the Hoffman bound satisfies
$\LHof \leq \bar{\vartheta} \leq \chr_f$, all three bounds---Hoffman,
Lov\'asz theta, and the fractional chromatic number---are \emph{permanently
shielded below~$4$} by the abundance of independent shadow vertices, even as
$\chr(M_k) = k \to \infty$.
\end{proof}

\begin{table}[ht]
\centering
\caption{Spectral shielding limits for Mycielski graphs.}
\label{tab:shielding}
\begin{tabular}{lcc}
\toprule
Bound & Limit for $M_k$ & Maximum certification \\
\midrule
$\LHof$            & $\approx 2.66$  & $\chr \geq 3$ \\
$\bar{\vartheta}$  & $\leq 3.377$    & $\chr \geq 3$ \\
$\chr_f$           & $\approx 3.377$ & $\chr \geq 3$ \\
$\chr$             & $k \to \infty$  & exact \\
\bottomrule
\end{tabular}
\end{table}

\noindent
\textbf{Consequence for Tihany.}
The Hoffman bound, Lov\'asz theta, and fractional chromatic number can certify
the symmetric $(3,3)$-partition because both parts only need $\chr \geq 3$,
which lies below the shielding ceiling.  For any asymmetric partition requiring
$\chr \geq 4$ on one side, these three bounds are provably insufficient on
Mycielski graphs.  Whether a $(2,k{-}1)$-Tihany partition exists for
general~$k$ remains open: vertex-criticality
(Theorem~\ref{thm:vertex-crit}) handles single-vertex removal, but removing an
edge is a strictly stronger condition
(Proposition~\ref{prop:asym-cond}).  For $k = 5$, exhaustive verification
confirms that every edge works (Observation~\ref{obs:edge-verify}).

\subsection{The Golden Sub-Induction (referee-safe version)}\label{ssec:sub-induction}

We now give a quantifier-tight induction for Theorem~\ref{thm:main}.

\begin{definition}[Canonical original-block embedding]
For $k\ge 4$, write $M_k=\myc(M_{k-1})$ and let
\[
V(M_k)=V_0(M_k)\sqcup V_1(M_k)\sqcup\{w_k\},
\]
where $V_0(M_k)$ is the original block (isomorphic to $M_{k-1}$),
$V_1(M_k)$ is the shadow block, and $w_k$ is the apex.
Let
\[
\iota_k:V(M_{k-1})\xrightarrow{\cong}V_0(M_k)
\]
denote this canonical isomorphism.
\end{definition}

\begin{lemma}[Coherent peeled family]\label{lem:coherent-family}
There exists a sequence $(P_k)_{k\ge 5}$ such that for every $k\ge 5$:
\begin{enumerate}[\upshape(i)]
\item $P_k\subseteq V(M_k)$ and $M_k[P_k]\cong C_5$;
\item for $k\ge 6$, $P_k\subseteq V_0(M_k)$ and
      $P_k=\iota_k(P_{k-1})$.
\end{enumerate}
\end{lemma}

\begin{proof}
Choose $P_5 = \{3,6,9,11,13\}$ (Lemma~\ref{lem:base-k5}). For $k\ge 6$,
define recursively
\[
P_k:=\iota_k(P_{k-1})\subseteq V_0(M_k).
\]
By construction $M_k[P_k]\cong M_{k-1}[P_{k-1}]\cong C_5$, and
$P_k \subseteq V_0(M_k)$ by definition.
\end{proof}

\begin{remark}[Golden peeling witness]\label{rem:golden-witness}
The steering theorem (Theorem~\ref{thm:steering}) provides evidence that each
$P_k$ can be realised by a golden peeling witness, i.e.\ there exists
$(z_k,\tau_k)$ with $z_k \in \mathcal{S}_k$ such that the threshold cut
$\{v : z_k(v) \geq \tau_k\} = P_k$.  The SO/OS control map is governed by
the same invertible matrix
$M = \bigl(\begin{smallmatrix}-\vp&\vp^{-1}\\ \vp^{-1}&-\vp\end{smallmatrix}\bigr)$
at every~$k$, and all subsequent positive-branch lifts apply identical scalar
attenuation to both paths, so the phase-layer targeting persists across levels.
This is verified computationally for $k = 5,\ldots,12$
(\cref{tab:verification}).  However, the proof of Theorem~\ref{thm:main}
requires only the combinatorial items (i)--(ii) of
Lemma~\ref{lem:coherent-family}, not the spectral witness.
\end{remark}

\begin{remark}[$1/\vp$ shadow attenuation]\label{rem:shadow-attenuation}
The coherence in Lemma~\ref{lem:coherent-family}(ii) can also be
seen directly from the eigenvector structure.  For $k \geq 6$, the last
Mycielski step uses the positive branch for both paths, so every basis vector
$b_j$ of~$\mathcal{S}_k$ satisfies
$(b_j)_{i+n_{k-1}} = (b_j)_i / \vp$ for all $i < n_{k-1}$.
Since $1/\vp < 1$, the shadow copy of any vertex is strictly less extreme
than the original.  The~$5$ extremal vertices of any $z \in \mathcal{S}_k$
therefore always lie in the original block.  At $k = 5$
(only~$2$ Mycielski steps), Path~B uses the negative branch at the last step,
giving shadow/original ratio $-\vp$ (amplification), so $k = 5$ requires
a separate base case.
\end{remark}

\begin{lemma}[Base case $k=5$, exact]\label{lem:base-k5}
Let $M_5$ be the $23$-vertex Mycielski graph (canonical numbering), and let
\[
P_5:=\{3,6,9,11,13\}.
\]
Then $M_5[P_5]\cong C_5$, and
\[
\chr(M_5\setminus P_5)=3.
\]
Hence $M_5$ satisfies the $(3,3)$-Tihany condition with partition
$S=P_5$, $T=V(M_5)\setminus P_5$.
\end{lemma}

\begin{proof}
The induced subgraph on $P_5$ has edges
\[
(3,9),\;(9,13),\;(13,6),\;(6,11),\;(11,3),
\]
so $M_5[P_5]\cong C_5$ and $\chr(M_5[P_5])=3$.

Set $T:=V(M_5)\setminus P_5$.  The cycle
\[
0\text{-}8\text{-}21\text{-}22\text{-}19\text{-}0
\]
is an odd cycle in $M_5[T]$, so $\chr(M_5[T])\ge 3$.

A proper $3$-coloring of $M_5[T]$ is given by color classes
\[
\begin{aligned}
C_1&=\{0,4,10,22\},\\
C_2&=\{1,12,14,15,17,19,20,21\},\\
C_3&=\{2,5,7,8,16,18\}.
\end{aligned}
\]
Thus $\chr(M_5[T])\le 3$, hence $\chr(M_5[T])=3$.
\end{proof}

\begin{lemma}[Containment along the coherent family]\label{lem:containment-coherent}
For $k\ge 6$, let $H_{k-1}:=M_{k-1}\setminus P_{k-1}$ and
$G_k:=M_k\setminus P_k$.  Then $G_k$ contains an induced copy of
$\myc(H_{k-1})$.  Consequently,
\[
\chr(G_k)\ge \chr\!\bigl(\myc(H_{k-1})\bigr)=\chr(H_{k-1})+1.
\]
\end{lemma}

\begin{proof}
By Lemma~\ref{lem:coherent-family}(ii),
$P_k = \iota_k(P_{k-1}) \subseteq V_0(M_k)$, so deleting~$P_k$ removes
exactly those original-block vertices corresponding to~$P_{k-1}$.  Write
$u_x$ for the shadow of vertex~$x$ and~$w_k$ for the apex in $M_k = \myc(M_{k-1})$.
Define $\psi_k \colon V(\myc(H_{k-1})) \to V(G_k)$ by
\[
  \psi_k(x) = \iota_k(x),
  \qquad
  \psi_k(u_x) = u_{\iota_k(x)},
  \qquad
  \psi_k(w) = w_k.
\]
This map is injective (its three branches target disjoint vertex sets:
original block, shadow block, and apex).
Edge preservation follows by cases:
\begin{itemize}
\item \emph{Original--original} ($xx' \in E(H_{k-1})$):
  $\iota_k(x)\iota_k(x') \in E(G_k)$ since $\iota_k$ is a graph isomorphism
  on the original block.
\item \emph{Shadow--original} ($u_x x' \in E(\myc(H_{k-1}))$, i.e.\
  $xx' \in E(H_{k-1})$): the Mycielski shadow rule gives
  $u_{\iota_k(x)}\iota_k(x') \in E(M_k)$; neither endpoint is in~$P_k$,
  so the edge survives in~$G_k$.
\item \emph{Apex--shadow} ($w u_x \in E(\myc(H_{k-1}))$):
  $w_k$ is adjacent to every shadow vertex in~$M_k$, hence
  $w_k u_{\iota_k(x)} \in E(G_k)$.
\item No shadow--shadow edges exist in either graph (by the Mycielski
  construction).
\end{itemize}
Moreover, the five ``orphan'' shadow vertices $u_{\iota_k(x)}$ for
$x \in P_{k-1}$ lie \emph{outside} $\psi_k(V(\myc(H_{k-1})))$, so they
cannot remove edges from the copy.  Hence $G_k[\psi_k(V(\myc(H_{k-1})))]
\cong \myc(H_{k-1})$---an \emph{induced} copy.

By chromatic monotonicity, $\chr(G_k) \geq \chr(\myc(H_{k-1}))$.  Since
$H_{k-1}$ has an edge (the base case $H_5 = M_5 \setminus P_5$ contains the
odd cycle $0$-$8$-$21$-$22$-$19$-$0$ by Lemma~\ref{lem:base-k5}, and adding
edges via Mycielski steps preserves this), the Mycielski chromatic
increment~\cite{mycielski1955} gives
$\chr(\myc(H_{k-1})) = \chr(H_{k-1}) + 1$.
\end{proof}

\subsubsection*{Proof of Theorem~\ref{thm:main}}

\begin{proof}[Proof of Theorem~\ref{thm:main}]
Let $(P_k)_{k\ge 5}$ be given by Lemma~\ref{lem:coherent-family}. Set
\[
S_k:=P_k,\qquad T_k:=V(M_k)\setminus P_k.
\]
Then $\chr(M_k[S_k])=\chr(C_5)=3$ for all $k\ge 5$.

It remains to show $\chr(M_k[T_k])\ge k-2$.

\medskip\noindent
\emph{Base case ($k=5$).}
This is Lemma~\ref{lem:base-k5}: $\chr(M_5\setminus P_5)=3=5-2$.

\medskip\noindent
\emph{Inductive step ($k \geq 6$).}
Assume $\chr(M_{k-1}\setminus P_{k-1})\ge (k-1)-2=k-3$.
By Lemma~\ref{lem:containment-coherent},
\[
\chr(M_k\setminus P_k)
\ge \chr\!\bigl(\myc(M_{k-1}\setminus P_{k-1})\bigr)
= \chr(M_{k-1}\setminus P_{k-1})+1
\ge (k-3)+1
= k-2.
\]

\medskip\noindent
\emph{Tihany verification.}
We have $\chr(M_k[S_k])\ge 3$ and $\chr(M_k[T_k])\ge k-2$.
Since $\chr(M_k)=k$ and $\om(M_k)=2<k=3+(k-2)-1$, this is a
$(3,k-2)$-Tihany partition for all $k\ge 5$.
\end{proof}

\begin{remark}\label{rem:prior-cases}
For $k \leq 7$, the pairs $(3,3),(3,4),(3,5)$ were already known from
Stiebitz~\cite{stiebitz1987} and Sachs~\cite{sachs1993}.  The pair $(3,6)$ at
$k = 8$ is, to our knowledge, the first new case settled by this method.  All
pairs $(3,t)$ for $t\geq 6$ are new.
\end{remark}

\begin{remark}\label{rem:edge-crit}
If one additionally proves that $\chr(M_k - e) = k - 1$ for every edge~$e$
of~$M_k$ (see Open Problem~2 in \cref{sec:open}), the same peeling gives
$(2,k{-}1)$-Tihany partitions: $S = \{u,v\}$ (an edge with $\chr = 2$),
$T = M_k\setminus\{u,v\}$ with $\chr \geq k-1$.
\end{remark}

%% ══════════════════════════════════════════════════════════════════════════
\section{The bigger picture: \texorpdfstring{$\vp$}{φ} as the irreducibility
threshold}\label{sec:bigger}
%% ══════════════════════════════════════════════════════════════════════════

\subsection{Four manifestations}

The golden ratio marks the boundary between reducible and irreducible structure
across domains; see \cref{tab:four}.

\begin{table}[ht]
\centering
\caption{Four manifestations of~$\vp$ as an irreducibility threshold.}
\label{tab:four}
\begin{tabular}{@{}p{2.8cm}p{3cm}p{3cm}p{4.2cm}@{}}
\toprule
Domain & Reducible & Irreducible & $\vp$ governs\ldots \\
\midrule
Graph colouring & $\chr = \om$ (perfect) & $\chr > \om$ (imperfect) &
  Eigenvalue of $C_5$: $\lambda_{\min} = -\vp$ \\
Tiling dynamics & Crystallographic & Icosahedral &
  Projection eigenvalue ($D_6 \to H_3$) \\
Information geometry & Generic $D_N$ families & $D_{12}$ (mod-$2 \cap$ mod-$3$) &
  Curvature minimum: $q^* = \vp^{-2}$ \\
Framing topology & $D \geq 4$: $\Z_2$ & $D = 3$: $\Z$ &
  Pentagonal symmetry of $H_3$ \\
\bottomrule
\end{tabular}
\end{table}

\subsection{The common root}

All four originate from
\[
  \cos\!\Bigl(\frac{2\pi}{5}\Bigr) = \frac{\vp - 1}{2}.
\]
Five-fold symmetry is the minimal structure where:
\begin{itemize}
\item Periodicity fails (crystallographic restriction).
\item Bipartiteness fails (odd cycle).
\item Parity and $3$-cycle coexist (mod-$2 \cap$ mod-$3$).
\item Framing is infinite ($\Z$ vs.\ $\Z_2$).
\end{itemize}

\subsection{The mod-2 / mod-3 bridge and Bruna's theorem}\label{ssec:bruna}

$C_5$ is where two arithmetic constraints first collide:
\begin{itemize}
\item \textbf{mod~2}: $\om = 2$ (triangle-free; bipartite obstruction).
\item \textbf{mod~3}: $\chr = 3$ (odd cycle; not $2$-colourable).
\end{itemize}
Bruna~\cite{bruna2025} proved that the Schur curvature on
$D_N$-equivariant exponential families has a unique minimum at
$q^* = \vp^{-2}$, with $D_{12}$ the minimal dihedral lattice where mod-$2$
(parity) and mod-$3$ (three-cycle) constraints coexist.

The pentagon~$C_5$ exhibits the same collision: $\om = 2$ (mod-$2$ constraint)
and $\chr = 3$ (mod-$3$ constraint).  This suggests a deeper connection between
the curvature landscape of dihedral exponential families and the chromatic
obstruction in pentagon-containing graphs.

\begin{remark}[Working hypothesis]\label{rem:bruna}
We conjecture that Bruna's variational result explains why the mod-$2$/mod-$3$
collision locks onto~$\vp$ in the graph-theoretic setting as well.  However,
this remains a \emph{heuristic bridge}---we do not yet have a formal transfer
theorem from the curvature framework to graph colouring.  The proof chain of
Theorem~\ref{thm:main}
(Lemma~\ref{lem:golden-propagation}, Theorem~\ref{thm:c5-peeling}, Lemma~\ref{lem:coherent-family}, and Lemma~\ref{lem:containment-coherent})
is entirely self-contained and does not depend on this connection.
\end{remark}

\subsection{Scope and limitations: why Mycielski, and why not Kneser}

The golden spectral method is precisely tailored to graph constructions that
\emph{recursively embed} $C_5$'s spectral structure.  The Mycielski functor is
the canonical such construction: its eigenvalue relation
$\mu^2 - \lambda\mu - \lambda^2 = 0$ is literally the defining equation
of~$\vp$, and its recursive block structure enables both the spectral
interferometry of Theorem~\ref{thm:steering} and the structural induction of
Theorem~\ref{thm:main}.

\emph{Kneser graphs} $K(n,r)$---another important family with
$\chr > \om$---illustrate the boundary of applicability.  Their spectra are
governed by the Johnson scheme, with eigenvalues
$(-1)^j\binom{n-r-j}{r-j}$ for $j = 0,\ldots,r$.  These are \emph{integers}
and generically do not lie in $\Q(\sqrt{5})$.

\begin{table}[ht]
\centering
\caption{Kneser graphs: integer spectra, no golden structure.}
\label{tab:kneser}
\begin{tabular}{lccll}
\toprule
Graph & $\chr$ & $\om$ & Eigenvalues & Golden? \\
\midrule
$K(5,2)$ (Petersen) & 3 & 2 & $\{3,1,-2\}$ & No ($\Z$-spectrum) \\
$K(7,3)$ & 3 & 2 & $\{10,2,-4\}$ & No \\
$K(8,3)$ & 4 & 2 & $\{20,4,-6\}$ & No \\
\bottomrule
\end{tabular}
\end{table}

The golden eigenvector machinery does not apply to Kneser graphs.  Their Tihany
partitions, if they exist, must arise from a different spectral witness or from
the combinatorial structure of set intersections.  This is not a weakness but a
\emph{design feature}: the golden spectral method solves the Tihany conjecture
where the chromatic obstruction is algebraically golden, and cleanly identifies
where it cannot reach.

\subsection{Future directions: \texorpdfstring{$\vp$}{φ}-recursive graph
operators}

The Mycielski functor acts on the adjacency matrix via a specific block
structure whose eigenvalue relation happens to generate~$\vp$.  Any graph
construction with a block form
\[
  \begin{pmatrix} A & \alpha A & \cdots \\ \beta A & O & \cdots \\ \vdots & & \ddots \end{pmatrix}
\]
that yields a characteristic equation $\mu^2 - \mu - 1 = 0$ (after rescaling)
will preserve golden eigenvalues and potentially support the same
peeling-and-induction strategy.  This includes:
\begin{itemize}
\item Generalised Mycielskians $M^{(p)}(G)$ (coning over~$p$ shadow copies).
\item Weighted Mycielskians with tuned edge weights.
\item Iterated cone constructions used in fractional chromatic number bounds.
\end{itemize}

\begin{conjecture}\label{conj:phi-recursive}
The class of $\vp$-recursive graphs---those generated by any
$\Q(\sqrt{5})$-preserving spectral operator starting from~$C_5$---satisfies
the Tihany conjecture.
\end{conjecture}

%% ══════════════════════════════════════════════════════════════════════════
\section{Open problems}\label{sec:open}
%% ══════════════════════════════════════════════════════════════════════════

\subsection{Fully proved in this paper}

\begin{enumerate}[(a)]
\item \emph{Golden propagation} (Lemma~\ref{lem:golden-propagation}): proven
  algebraically (\cref{sec:propagation}).
\item \emph{$C_5$-peeling existence}: proven constructively via the Steering
  Theorem (Theorem~\ref{thm:steering}, $\det = \sqrt{5}$).
\item \emph{Coherent peeled family} (Lemma~\ref{lem:coherent-family}): proven
  for all $k \geq 5$, with original-block localisation for $k \geq 6$.
\item \emph{Containment along coherent family}
  (Lemma~\ref{lem:containment-coherent}): proven for all $k \geq 6$.
\item \emph{$(3,k{-}2)$-Tihany for all~$M_k$}: Theorem~\ref{thm:main}, via
  sub-induction with base case $k = 5$ (direct computation).
\item \emph{Spectral shielding}: $\chr_f(M_k) \to 3.377 < 4$; spectral
  methods provably cannot certify $\chr \geq 4$ (Proposition~\ref{prop:shielding}).
\end{enumerate}

\subsection{Computationally verified (proofs incomplete)}

\begin{enumerate}[(a)]
\item \emph{Hoffman margin $F = \vp^{-3}$} (Observation~\ref{obs:margin}): verified for
  $k = 5,\ldots,12$; analytical proof would give a purely algebraic spectral
  certificate for the $(3,3)$-partition.
\item \emph{Hoffman bound of large part $> \sqrt{5}$}
  (Observation~\ref{obs:hoffman-large}, Conjecture~\ref{conj:hoffman-large}): verified for $k \leq 12$; the asymptotic
  convergence argument (\cref{ssec:spectral-growth}) is strong but the
  base-case bound needs formalisation.
\end{enumerate}

\subsection{Remaining open problems}

\begin{enumerate}[(1)]
\item \textbf{(Resolved.)} \textbf{All Tihany pairs for $M_k$.}
  The companion paper~\cite{santos2026amplification} proves that
  $M_k$ satisfies the Tihany conjecture for \emph{every} pair $(s,t)$
  with $s + t = k + 1$ and $s, t \geq 2$, by a purely structural argument:
  the canonical embedding $M_s \hookrightarrow M_k$ (iterated original block)
  gives $\chr(M_k[S]) = s$, and the Mycielski complement
  $R_{k,s} = M_k \setminus M_s$ satisfies
  $\chr(R_{k,s}) \geq k - s + 1 = t$ via the containment
  $R_{k,s} \supseteq M(R_{k-1,s})$.  This supersedes
  Theorem~\ref{thm:main}, which covers only the $(3, k{-}2)$ family.
  The $(2, k{-}1)$ pair, listed below as Open Problem~2, is now
  an immediate corollary.

\item \textbf{Analytical proof of $F = \vp^{-3}$ margin.}
  Verified to $k = 12$; the 4D Phase-Lock is proven but the quantitative bound
  needs formalisation.  This would upgrade Observation~\ref{obs:margin} from
  computational observation to theorem.

\item \textbf{Edge-criticality: $\chr(M_k - e) = k - 1$ for all edges.}
  Verified for $k = 5$.  This is now subsumed by the structural
  partition approach of~\cite{santos2026amplification}, which gives
  $(2, k{-}1)$-Tihany partitions directly.  Edge-criticality remains
  interesting as an independent structural property of Mycielski graphs.
\end{enumerate}

\subsection{Future work (beyond this paper)}

\begin{itemize}
\item \textbf{$\vp$-recursive graph operators.}
  Define and classify all graph constructions whose eigenvalue relation
  generates~$\vp$.  Prove Tihany for this broader class
  (Conjecture~\ref{conj:phi-recursive}).

\item \textbf{Non-golden families (Kneser, Schrijver).}
  These have integer spectra and require non-golden spectral witnesses or
  topological methods.

\item \textbf{Bruna transfer theorem.}
  Formalise the heuristic bridge between dihedral curvature minima and
  pentagonal chromatic obstructions (Remark~\ref{rem:bruna}).

\item \textbf{The full Erd\H{o}s--Lov\'asz conjecture.}
  For general $\chr > \om$ graphs, a fundamentally different approach may be
  needed.  The Strong Perfect Graph Theorem guarantees an odd hole or odd
  antihole; whether its spectral trace always suffices for Tihany remains
  unknown.
\end{itemize}

%% ══════════════════════════════════════════════════════════════════════════
%% References
%% ══════════════════════════════════════════════════════════════════════════
\bibliographystyle{amsplain}
\bibliography{references}

%% ══════════════════════════════════════════════════════════════════════════
\appendix
\section{Reproducibility appendix}\label{app:repro}
%% ══════════════════════════════════════════════════════════════════════════

All computational claims in this paper are verified by the scripts in the
accompanying \texttt{code/} directory.  This appendix documents the exact
protocol, dependencies, and expected output so that every numerical statement
can be independently reproduced.

\subsection{Environment}

\begin{tabular}{@{}ll@{}}
\toprule
\textbf{Component} & \textbf{Version / specification} \\
\midrule
Python      & $\geq 3.10$ \\
NumPy       & $\geq 1.24$ \\
SciPy       & $\geq 1.10$ \\
NetworkX    & $\geq 3.0$ \\
\bottomrule
\end{tabular}

\medskip\noindent
Install dependencies: \texttt{pip install -r code/requirements.txt}.

\subsection{Scripts and their claims}

\begin{enumerate}[\bfseries S1.]
\item \texttt{c5\_peeling\_verification.py} --- verifies Theorem~\ref{thm:c5-peeling}
  and Observation~\ref{obs:margin} for $k = 3,\ldots,12$.

  \textbf{Protocol.}  For each Mycielski graph~$M_k$:
  \begin{enumerate}[(a)]
  \item Construct $M_k$ via \texttt{networkx.mycielski\_graph(k)}.
  \item Compute the adjacency eigendecomposition (NumPy \texttt{eigh}).
  \item Identify the $+\vp^{k-4}$ eigenspace (tolerance $<0.05$).
  \item Sample random unit vectors in the eigenspace
    (\texttt{RandomState(seed=42)}, up to $5000$~trials).
  \item For each sample, extract the bottom-$5$ (or top-$5$) vertices by
    eigenvector coordinate and test whether the induced subgraph is
    isomorphic to~$C_5$ (\texttt{nx.is\_isomorphic}).
  \item On success, compute $\LHof$ of the remainder and record the
    margin~$F$.
  \end{enumerate}

  \textbf{Expected output.}  For all $k = 5,\ldots,12$:
  $F(\tau^*) = 0.2361 \approx \vp^{-3}$ (to four decimal places) and the
  peeled subgraph is $C_5$.

  \textbf{Invocation:}
  \begin{quote}\ttfamily
  python code/c5\_peeling\_verification.py 12
  \end{quote}
  Runtime: $\approx 2$~minutes on a standard laptop (the $k = 12$ graph has
  $3071$~vertices).

\item \texttt{explore\_spectral\_tihany.py} --- computes eigenvalues, Hoffman
  bounds, Lov\'asz theta approximations, and golden-ratio proximity for a
  library of test graphs ($C_5$, Petersen, Gr\"otzsch, Mycielski, Kneser).

  \textbf{Invocation:}
  \begin{quote}\ttfamily
  python code/explore\_spectral\_tihany.py
  \end{quote}

\item \texttt{test\_spectral\_cuts.py} --- tests spectral threshold cuts
  across the graph library.

\item \texttt{lemma3\_analysis.py} --- verifies the spectral certificate of
  Lemma~\ref{lem:spectral-certificate} on small cases.
\end{enumerate}

\subsection{Deterministic seed and exact reproduction}

All randomised searches use the fixed seed \texttt{seed\,=\,42}
(\texttt{numpy.random.RandomState(42)}).  With identical NumPy and NetworkX
versions, the trial indices in \cref{tab:verification} are reproduced
\emph{exactly}.  Minor version differences in LAPACK may permute eigenvectors
within a degenerate eigenspace, potentially changing trial indices, but the
\emph{existence} result (that a $C_5$-peeling direction is found within
$5000$~trials) is robust.

\subsection{Base-case verification ($k = 5$)}

The base case $k = 5$ of Theorem~\ref{thm:main} requires
$\chr(M_5 \setminus P_5) \geq 3$.  Lemma~\ref{lem:base-k5} provides a
self-contained paper certificate: the odd cycle $0$-$8$-$21$-$22$-$19$-$0$
proves $\chr \geq 3$, and the explicit $3$-coloring proves $\chr \leq 3$.
Script~\textbf{S1} independently confirms $\LHof(M_5 \setminus P_5) = 2.2855 > 2$
(a spectral lower bound), and verifies that the $18$-vertex remainder is not
bipartite (ruling out $\chr = 2$).  Note: \texttt{greedy\_color} provides an
upper bound only, not an optimality certificate; the exact value $\chr = 3$ is
established by Lemma~\ref{lem:base-k5}.

\subsection{Edge-criticality verification ($k = 5$)}

For all $71$ edges of~$M_5$, we verified
$\chr(M_5 - \{u,v\}) = 4 = \chr(M_5) - 1$ by exhaustive
$3$-colouring search.  This supports the edge-removal property of
Observation~\ref{obs:edge-verify}.

\end{document}
